% =========================================================================
\section{Which degrees give a finite construction}\label{sec:degrees}

Theorems~\ref{thm:arc} and~\ref{thm:mean} are statements about real functions of a real parameter $p$.
A drawing needs more: the residual $t=X/u^{p}$ must be \emph{constructed} at every step. The triangle
stage constructs integer powers only. This section says exactly what follows.

\begin{proposition}[Rational degrees are finite constructions]\label{prop:rational}
Let $p=\mu/\nu$ with coprime positive integers $\mu,\nu$, and $X>0$. Then $X^{1/p}=(X^{\nu})^{1/\mu}$,
and:
\begin{enumerate}[label=(\roman*)]
\item if $\mu\ge3$, every integer-degree protocol of this paper applies with degree $\mu$ and input
  $X^{\nu}$. In particular the decentered and three-mean iterations converge monotonically from every
  positive start, with orders five and nine, at no more than $10\mu+9$ and $10\mu+11$ traces per step
  (Section~\ref{sec:work}). The input $X^{\nu}$ is fixed preparation: $\nu-1$ five-join products, once;
\item if $\mu\in\{1,2\}$, the root $X^{\nu}$ or $\sqrt{X^{\nu}}$ is constructed exactly, by $\nu-1$
  products and at most one mean proportional.
\end{enumerate}
For $p<0$ apply the above to $|p|$ and invert the output (five joins, once).
\end{proposition}
\begin{proof}
The residual of the degree-$\mu$ problem with input $X^{\nu}$ is $t=X^{\nu}s^{\mu}$, built by the
triangle stage with $\mu$ in place of $p$ from the prepared point $M=(0,H(1-X^{-\nu}))$. Theorems
\ref{thm:arc} and \ref{thm:mean} apply with real degree $\mu>2$ and positive input $X^{\nu}$, and their
fixed point is $(X^{\nu})^{1/\mu}$.
\end{proof}

Two consequences are worth stating. First, the degree seen by the iteration is $\mu$, not $p$: the
degree $p=3/2$ is handled as a cube root of $X^{2}$, inside the domain $\mu>2$ of the unmodified
formulas although $3/2<2$. Second, the cost is governed by $\mu$: $p=355/113$ costs as much per step
as degree $355$.

The restriction to degree greater than two is real for the \emph{raw} formulas. At $p=2$ the
coefficients $m$ and $n$ of \eqref{eq:mcoeff} vanish, $g=1$, and $q_M=\sqrt t$ is the exact support,
while the preparation of the decentered arc is singular. For $\tfrac12<p<2$ the weights $m$ and $n$ are
negative and the positivity argument of Theorem~\ref{thm:mean} no longer applies. As a scalar statement
the restriction can be removed:

\begin{proposition}[Scalar extension to every nonzero real degree]\label{prop:scalar}
For every real $p\neq0$, every $X>0$ and every start $u_0>0$ there is an explicit modification of the
three-mean iteration \eqref{eq:qM} that converges monotonically to $X^{1/p}$ without crossing it.
\end{proposition}
\begin{proof}
Choose an integer $k\ge0$ with $p'=2^{k}|p|>2$, and write $F_{p',X}(v)=1/\Tm_M(1/v)$ for the three-mean
iteration of degree $p'$ in the direct state. For $p>0$ put $v_0=u_0^{1/2^{k}}$, iterate
$v_{n+1}=F_{p',X}(v_n)$ and report $u_n=v_n^{2^{k}}$; for $p<0$ put $v_0=u_0^{-1/2^{k}}$ and report
$u_n=v_n^{-2^{k}}$. The reported start is exactly $u_0$. Theorem~\ref{thm:mean} gives
$v_n\to X^{1/p'}$ monotonically, so $u_n\to X^{\pm2^{k}/p'}=X^{1/p}$. Positive integer powers preserve
inequalities and inversion reverses them together with the direction of approach. For $p>2$ take
$k=0$. The same argument applies to the decentered arc with Theorem~\ref{thm:arc}.
\end{proof}

Proposition~\ref{prop:scalar} is a change of variable and adds no geometric content beyond
Proposition~\ref{prop:rational}: for irrational $p$ the step still needs the real power
$v^{p'}$, which is a problem of the same nature as the root being sought. Irrational degrees can
nevertheless be reached by the fast iterations through a rational approximation of $p$, with a
controlled error:

\begin{proposition}[Irrational degrees by rational approximation]\label{prop:irrational}
Let $p>0$ be real, $X>0$, $r=X^{1/p}$, and let $\mu,\nu$ be positive integers with
$\bigl|\nu/\mu-1/p\bigr|\le\eta$. Then $r'=X^{\nu/\mu}$ satisfies
\[ \Bigl|\frac{r'}{r}-1\Bigr|\le e^{\eta|\log X|}-1 . \]
If moreover $\tilde u>0$ satisfies $|\tilde u/r'-1|\le\eps_1$, then
$|\tilde u/r-1|\le(1+\eps_1)\,e^{\eta|\log X|}-1$.
\end{proposition}
\begin{proof}
$r'/r=X^{\nu/\mu-1/p}$, so $|\log(r'/r)|\le\eta|\log X|$, which gives the first bound since
$1-e^{-y}\le e^{y}-1$ for $y\ge0$. For the second, $\tilde u/r=(\tilde u/r')(r'/r)$ lies between
$(1-\eps_1)e^{-y}$ and $(1+\eps_1)e^{y}$ with $y=\eta|\log X|$, and
$1-(1-\eps_1)e^{-y}\le(1+\eps_1)e^{y}-1$.
\end{proof}

\begin{remark}[Cost of the rational route]\label{rem:pi}
Take $p=\pi$, $X=2$, $r=2^{1/\pi}=1.2468689889\ldots$ With $\mu/\nu=22/7$, $\eta\approx1.28\times10^{-4}$
and the approximation error is at most $8.9\times10^{-5}$; a step costs at most $229$ traces (decentered)
or $231$ (three means), after $30$ joins of preparation. With $\mu/\nu=355/113$,
$\eta\approx2.70\times10^{-8}$ and the error is at most $1.9\times10^{-8}$, but a step costs at most
$3559$ or $3561$ traces after $560$ joins of preparation. The direct subdivision of
Section~\ref{sec:enclosure} reaches relative width $10^{-8}$ for the true exponent in at most $123$
traces. Since a convergent $\nu/\mu$ of $1/p$ has $\eta$ of order $\mu^{-2}$, the rational route needs
$\mu$ of order $\eps^{-1/2}$ for accuracy $\eps$, while its step cost is linear in $\mu$. For irrational
degrees the subdivision is therefore the recommended calculator, and the rational route is of interest
only when a small $\mu$ is acceptable.
\end{remark}

A third route would approximate the real power $u^{p}$ inside each step by an inner subdivision. That
is a two-level recursion whose outer iteration receives an inexact residual; its convergence would need
a perturbation analysis that we have not carried out (Section~\ref{sec:conclusion}).
